% llm_peer/drafts/en/manuscript.tex
% P4.2 manuscript skeleton (plan-r3 derivative-paper extension, 2026-05-15)
% Compile: xelatex + biber + xelatex × 2

% papers_assets/_shared/preamble_en.tex
% Shared English LaTeX preamble for all 5 paper drafts.
% P4.1 (plan-r3 derivative-paper extension, 2026-05-15)
% Compile: xelatex + biber + xelatex × 2

\documentclass[12pt,a4paper]{article}

% --- Fonts & Encoding ---
\usepackage{fontspec}
\setmainfont{Times New Roman}[Ligatures=TeX]

% --- Math ---
\usepackage{amsmath, amssymb, amsthm}

% --- Tables & Figures ---
\usepackage{booktabs}
\usepackage{graphicx}
\usepackage{longtable}
\usepackage{rotating}
\usepackage{caption}
\captionsetup{font=small,labelfont=bf,skip=4pt}

% --- Layout ---
\usepackage[margin=1in]{geometry}
\usepackage{setspace}
\onehalfspacing
\usepackage{hyperref}
\hypersetup{colorlinks=true,linkcolor=blue,citecolor=blue,urlcolor=blue}

% --- Bibliography (biblatex + biber) ---
\usepackage[backend=biber,style=apa,natbib=true]{biblatex}

% --- Useful macros ---
\newcommand{\code}[1]{\texttt{#1}}
\newcommand{\NoteThis}[1]{\textbf{\textcolor{red}{[TODO: #1]}}}
\usepackage[dvipsnames]{xcolor}

% --- Theorem environments ---
\newtheorem{prop}{Proposition}
\newtheorem{hyp}{Hypothesis}


\addbibresource{refs.bib}

\title{LLM-Generated Peer Identification and Causal Inference for Bank FinTech Strategy}
\author{\NoteThis{Author names}}
\date{\today}

\begin{document}
\maketitle

\begin{abstract}
\NoteThis{Final abstract pending.}
Using bank-year media coverage of FinTech and LLM-defined peer set (Cao \emph{et~al.}\ 2025) + leave-one-out instrumental variable, we measure
LLM-peer leave-one-out tone IV (\code{llm\_peer\_tone\_iv\_full}) for Chinese listed banks. We find \NoteThis{key result}\@.
\end{abstract}

\section{Introduction}\label{sec:intro}
\NoteThis{Motivate the research question and contribution.}

\section{Literature Review}\label{sec:lit}
(1)~LLM-based peer identification (Cao \emph{et~al.}\ 2025 \emph{RAS}); (2)~Stock--Yogo (2005) weak-instrument testing; (3)~Kleibergen--Paap (2006) rank-based identification.

\section{Methodology}\label{sec:method}
\subsection{Sample and Data}
The bank--year panel covers approximately 164 bank--years
(2017--2018). All bank--year observations are processed by the
\code{news\_process} pipeline (Phase 9b baseline; see
\code{papers\_assets/V6\_METHODOLOGY\_CHANGES.md}).

\subsection{Tone Annotation (F1: Three-Family LLM Validation)}
We employ three independent LLM families for sentiment annotation:
\code{gpt-5.4-nano} (OpenAI), \code{gemini-3.1-flash-lite-preview}
(Google), and \code{deepseek-v4-flash} (DeepSeek\,/\,Grok). All three
routed through yunwu.ai (\code{/v1/chat/completions}, OpenAI-compatible),
\code{temperature=0}, \code{seed=42}. Inter-family agreement reported as
Cohen's $\kappa$ and Krippendorff's $\alpha$; majority vote with ties marked
\code{undefined} (not \code{neutral}).

\subsection{Article-Level Coding (F2: Decontamination)}
Articles where all three families fail or split three ways are coded
\emph{missing}, not \code{neutral}. Bank--FinTech relevance failing
majority verification is also missing. This eliminates the 81.7\,\%
contamination in pre-v6 ``neutral'' labels.

\subsection{Zero-Observation Semantics (F3: NaN Main Path)}
For bank--years with \code{n\_articles=0}, similarity-based features
(\code{auth\_sim\_*}, \code{benefit\_*}, \code{article\_sim\_*}),
ratios (\code{token\_adopt\_*}, \code{hidden\_adopt\_*},
\code{subst\_adopt\_*}), and \code{fintech\_intensity\_mean} are coded
as \code{NaN}. Counts remain at 0. Listwise deletion follows Allison
(2001) / Heckman (1979).

\subsection{Bartik IV Pruning (F4)}
Of approximately 30 candidate Bartik instruments, 28 are removed in v6
for violating Goldsmith-Pinkham \emph{et~al.}\ (2020, \emph{AER})
exogeneity. Retained: \code{bartik\_type\_iv\_full} and
\code{prov\_leave1out\_tone\_full}.

\section{Data}\label{sec:data}
\NoteThis{Descriptive statistics in Table 1 (see \code{papers/llm_peer/tables.py:build\_table1}).}

\section{Main Results}\label{sec:main}
\NoteThis{Table 2 from \code{output/tables/table2\_*.csv}; Figure 2 coefplot.}

\section{Mechanism}\label{sec:mech}
\NoteThis{Table 3 from \code{output/tables/table3\_*.csv}; Figure 3.}

\section{Heterogeneity}\label{sec:het}
\NoteThis{Table 4 (SOE × paper-specific interactions) from
\code{output/tables/table4\_heterogeneity.csv}.}

\section{Robustness}\label{sec:robust}
\NoteThis{Table 5 (paper-specific robustness battery) from
\code{output/tables/table5\_*.csv}.}

\section{Conclusion}\label{sec:concl}
\NoteThis{Final discussion and policy implications.}

\printbibliography

\end{document}
